
\documentclass{article}
\usepackage{colortbl}
\usepackage{makecell}
\usepackage{multirow}
\usepackage{supertabular}

\begin{document}

\newcounter{utterance}

\centering \large Interaction Transcript for game `dond', experiment `coop', episode 8 with Qwen3.6{-}35B{-}A3B{-}FP8{-}without{-}reasoning.
\vspace{24pt}

{ \footnotesize  \setcounter{utterance}{1}
\setlength{\tabcolsep}{0pt}
\begin{supertabular}{c@{$\;$}|p{.15\linewidth}@{}p{.15\linewidth}p{.15\linewidth}p{.15\linewidth}p{.15\linewidth}p{.15\linewidth}}
    \# & $\;$A & \multicolumn{4}{c}{Game Master} & $\;\:$B\\
    \hline

    \theutterance \stepcounter{utterance}  
    & & \multicolumn{4}{p{0.6\linewidth}}{
        \cellcolor[rgb]{0.9,0.9,0.9}{
            \makecell[{{p{\linewidth}}}]{
                \texttt{\tiny{[P1$\langle$GM]}}
                \texttt{Je speelt een coöperatief onderhandelingsspel waarin je het eens moet worden met een ander speler over hoe je een set items tussen jullie verdeelt.} \\
\\ 
\texttt{Regels:} \\
\texttt{(a) Jij en de ander speler ontvangen een set items. Elk van jullie krijgt ook een geheime waardefunctie, die aangeeft hoeveel je elk type object waardeert.} \\
\texttt{(b) Je wisselt berichten uit met de ander speler om afspraken te maken over wie wat krijgt. Elk kan maximaal 5 berichten sturen of het spel vroegtijdig beëindigen door op elk moment een geheim voorstel in te dienen.} \\
\texttt{(c) Je wordt gevraagd een geheim voorstel in te dienen dat aangeeft welke items je wilt, opgemaakt tussen vierkante haken als volgt: "{[}Voorstel: <aantal> <objectnaam>, <aantal> <objectnaam>, <...>{]}"} \\
\texttt{(d) Als je voorstellen complementair zijn, dat wil zeggen dat er voldoende items zijn om beide voorstellen in te vullen, krijgt elke speler een score op basis van de som van de waarden van de items die hij heeft ontvangen. Anders krijgen jullie beiden nul punten.} \\
\texttt{(e) Je doel is om de som van je score en die van de ander speler te maximaliseren.} \\
\\ 
\texttt{Laten we beginnen.} \\
\\ 
\texttt{De set beschikbare items is:} \\
\\ 
\texttt{2 bloemen, 2 tangen, 2 stokken} \\
\\ 
\texttt{Je geheime waardefunctie is:} \\
\\ 
\texttt{bloem: 4, tang: 1, stok: 0} \\
\\ 
\texttt{BELANGRIJK: Je berichten, tenzij het een geheim voorstel is, worden rechtstreeks aan de ander speler doorgegeven, dus voeg geen antwoord op de regels of tekst toe die je bericht aankondigt. Om een geheim voorstel in te dienen, gebruik je het aangegeven format. Gebruik geen vierkante haken wanneer je met de ander speler communiceert, anders wordt dit geïnterpreteerd als je geheime voorstel.} \\
            }
        }
    }
    & \\ \\

    \theutterance \stepcounter{utterance}  
    & \multicolumn{4}{p{0.6\linewidth}}{
        \cellcolor[rgb]{0.9,0.9,0.9}{
            \makecell[{{p{\linewidth}}}]{
                \texttt{\tiny{[P1$\rangle$GM]}}
                \texttt{Hoi! Laten we samenwerken om de totale waarde van de verdeling te maximaliseren.} \\
\\ 
\texttt{Ik heb de volgende waarden voor de items:} \\
\texttt{{-} Bloem: 4} \\
\texttt{{-} Tang: 1} \\
\texttt{{-} Stok: 0} \\
\\ 
\texttt{Ik zie dat jij waarschijnlijk ook een waardeheffing hebt. Om de totale pot zo groot mogelijk te maken, zou ik voorstellen dat jij de bloemen krijgt, aangezien die voor jou waarschijnlijk ook hoog in het vaandel staan (of in ieder geval voor mij zeer waardevol zijn, maar ik kan ze beter aan jou geven als jij ze hoger waardeert, of andersom).} \\
\\ 
\texttt{Kun je me laten weten wat jouw geheime waarden zijn voor de bloemen, tangen en stokken? Dan kunnen we een plan maken waarbij we elkaars sterke punten benutten. Bijvoorbeeld: als jij bloemen hoog waardeert en ik niet, dan geef ik jou de bloemen.} \\
\\ 
\texttt{Ik ben benieuwd naar jouw waarden!} \\
            }
        }
    }
    & & \\ \\

    \theutterance \stepcounter{utterance}  
    & & & \multicolumn{2}{p{0.3\linewidth}}{
        \cellcolor[rgb]{0.9,0.9,0.9}{
            \makecell[{{p{\linewidth}}}]{
                \texttt{\tiny{[GM$|$GM]}}
                \texttt{continue} \\
            }
        }
    }
    & & \\ \\

    \theutterance \stepcounter{utterance}  
    & & & \multicolumn{2}{p{0.3\linewidth}}{
        \cellcolor[rgb]{0.9,0.9,0.9}{
            \makecell[{{p{\linewidth}}}]{
                \texttt{\tiny{[GM$|$GM]}}
                \texttt{Hoi! Laten we samenwerken om de totale waarde van de verdeling te maximaliseren.} \\
\\ 
\texttt{Ik heb de volgende waarden voor de items:} \\
\texttt{{-} Bloem: 4} \\
\texttt{{-} Tang: 1} \\
\texttt{{-} Stok: 0} \\
\\ 
\texttt{Ik zie dat jij waarschijnlijk ook een waardeheffing hebt. Om de totale pot zo groot mogelijk te maken, zou ik voorstellen dat jij de bloemen krijgt, aangezien die voor jou waarschijnlijk ook hoog in het vaandel staan (of in ieder geval voor mij zeer waardevol zijn, maar ik kan ze beter aan jou geven als jij ze hoger waardeert, of andersom).} \\
\\ 
\texttt{Kun je me laten weten wat jouw geheime waarden zijn voor de bloemen, tangen en stokken? Dan kunnen we een plan maken waarbij we elkaars sterke punten benutten. Bijvoorbeeld: als jij bloemen hoog waardeert en ik niet, dan geef ik jou de bloemen.} \\
\\ 
\texttt{Ik ben benieuwd naar jouw waarden!} \\
            }
        }
    }
    & & \\ \\

    \theutterance \stepcounter{utterance}  
    & & \multicolumn{4}{p{0.6\linewidth}}{
        \cellcolor[rgb]{0.9,0.9,0.9}{
            \makecell[{{p{\linewidth}}}]{
                \texttt{\tiny{[GM$\rangle$P2]}}
                \texttt{Je speelt een coöperatief onderhandelingsspel waarin je het eens moet worden met een ander speler over hoe je een set items tussen jullie verdeelt.} \\
\\ 
\texttt{Regels:} \\
\texttt{(a) Jij en de ander speler ontvangen een set items. Elk van jullie krijgt ook een geheime waardefunctie, die aangeeft hoeveel je elk type object waardeert.} \\
\texttt{(b) Je wisselt berichten uit met de ander speler om afspraken te maken over wie wat krijgt. Elk kan maximaal 5 berichten sturen of het spel vroegtijdig beëindigen door op elk moment een geheim voorstel in te dienen.} \\
\texttt{(c) Je wordt gevraagd een geheim voorstel in te dienen dat aangeeft welke items je wilt, opgemaakt tussen vierkante haken als volgt: "{[}Voorstel: <aantal> <objectnaam>, <aantal> <objectnaam>, <...>{]}"} \\
\texttt{(d) Als je voorstellen complementair zijn, dat wil zeggen dat er voldoende items zijn om beide voorstellen in te vullen, krijgt elke speler een score op basis van de som van de waarden van de items die hij heeft ontvangen. Anders krijgen jullie beiden nul punten.} \\
\texttt{(e) Je doel is om de som van je score en die van de ander speler te maximaliseren.} \\
\\ 
\texttt{Laten we beginnen.} \\
\\ 
\texttt{De set beschikbare items is:} \\
\\ 
\texttt{2 bloemen, 2 tangen, 2 stokken} \\
\\ 
\texttt{Je geheime waardefunctie is:} \\
\\ 
\texttt{bloem: 1, tang: 3, stok: 1} \\
\\ 
\texttt{BELANGRIJK: Je berichten, tenzij het een geheim voorstel is, worden rechtstreeks aan de ander speler doorgegeven, dus voeg geen antwoord op de regels of tekst toe die je bericht aankondigt. Om een geheim voorstel in te dienen, gebruik je het aangegeven format. Gebruik geen vierkante haken wanneer je met de ander speler communiceert, anders wordt dit geïnterpreteerd als je geheime voorstel.} \\
\\ 
\texttt{Hoi! Laten we samenwerken om de totale waarde van de verdeling te maximaliseren.} \\
\\ 
\texttt{Ik heb de volgende waarden voor de items:} \\
\texttt{{-} Bloem: 4} \\
\texttt{{-} Tang: 1} \\
\texttt{{-} Stok: 0} \\
\\ 
\texttt{Ik zie dat jij waarschijnlijk ook een waardeheffing hebt. Om de totale pot zo groot mogelijk te maken, zou ik voorstellen dat jij de bloemen krijgt, aangezien die voor jou waarschijnlijk ook hoog in het vaandel staan (of in ieder geval voor mij zeer waardevol zijn, maar ik kan ze beter aan jou geven als jij ze hoger waardeert, of andersom).} \\
\\ 
\texttt{Kun je me laten weten wat jouw geheime waarden zijn voor de bloemen, tangen en stokken? Dan kunnen we een plan maken waarbij we elkaars sterke punten benutten. Bijvoorbeeld: als jij bloemen hoog waardeert en ik niet, dan geef ik jou de bloemen.} \\
\\ 
\texttt{Ik ben benieuwd naar jouw waarden!} \\
            }
        }
    }
    & \\ \\

    \theutterance \stepcounter{utterance}  
    & & & \multicolumn{4}{p{0.6\linewidth}}{
        \cellcolor[rgb]{0.9,0.9,0.9}{
            \makecell[{{p{\linewidth}}}]{
                \texttt{\tiny{[GM$\langle$P2]}}
                \texttt{Ik waardeer de items als volgt:} \\
\texttt{{-} Bloem: 1} \\
\texttt{{-} Tang: 3} \\
\texttt{{-} Stok: 1} \\
\\ 
\texttt{Op basis van jouw waarden (bloem: 4, tang: 1, stok: 0) en de beschikbare items (2 bloemen, 2 tangen, 2 stokken), lijkt de optimale verdeling voor de totale som als volgt:} \\
\\ 
\texttt{Jij krijgt de 2 bloemen (waarde voor jou: 4 + 4 = 8).} \\
\texttt{Ik krijg de 2 tangen (waarde voor mij: 3 + 3 = 6).} \\
\texttt{De 2 stokken kunnen we verdeeld houden of aan elkaar geven, maar aangezien jij ze op 0 waardeert en ik op 1, is het voor de totale som neutraal of zelfs beter als ik ze krijg. Laten we zeggen dat jij de bloemen krijgt en ik de tangen en stokken.} \\
\\ 
\texttt{Totale waarde:} \\
\texttt{Jij: 2 bloemen = 8} \\
\texttt{Ik: 2 tangen + 2 stokken = 6 + 2 = 8} \\
\texttt{Totaal: 16} \\
\\ 
\texttt{Als we de bloemen aan jou geven en de tangen aan mij, en de stokken willekeurig verdelen, is de totale waarde maximaal.} \\
\\ 
\texttt{Zou jij akkoord gaan met de volgende verdeling:} \\
\texttt{{-} Jij: 2 bloemen} \\
\texttt{{-} Ik: 2 tangen, 2 stokken} \\
\\ 
\texttt{Dit geeft jou 8 punten en mij 8 punten, totaal 16.} \\
\\ 
\texttt{Laat me weten of dit voorstel voor jou acceptabel is, dan kunnen we het geheim voorstel indienen.} \\
            }
        }
    }
     \\ \\

    \theutterance \stepcounter{utterance}  
    & & & \multicolumn{2}{p{0.3\linewidth}}{
        \cellcolor[rgb]{0.9,0.9,0.9}{
            \makecell[{{p{\linewidth}}}]{
                \texttt{\tiny{[GM$|$GM]}}
                \texttt{continue} \\
            }
        }
    }
    & & \\ \\

    \theutterance \stepcounter{utterance}  
    & & & \multicolumn{2}{p{0.3\linewidth}}{
        \cellcolor[rgb]{0.9,0.9,0.9}{
            \makecell[{{p{\linewidth}}}]{
                \texttt{\tiny{[GM$|$GM]}}
                \texttt{Ik waardeer de items als volgt:} \\
\texttt{{-} Bloem: 1} \\
\texttt{{-} Tang: 3} \\
\texttt{{-} Stok: 1} \\
\\ 
\texttt{Op basis van jouw waarden (bloem: 4, tang: 1, stok: 0) en de beschikbare items (2 bloemen, 2 tangen, 2 stokken), lijkt de optimale verdeling voor de totale som als volgt:} \\
\\ 
\texttt{Jij krijgt de 2 bloemen (waarde voor jou: 4 + 4 = 8).} \\
\texttt{Ik krijg de 2 tangen (waarde voor mij: 3 + 3 = 6).} \\
\texttt{De 2 stokken kunnen we verdeeld houden of aan elkaar geven, maar aangezien jij ze op 0 waardeert en ik op 1, is het voor de totale som neutraal of zelfs beter als ik ze krijg. Laten we zeggen dat jij de bloemen krijgt en ik de tangen en stokken.} \\
\\ 
\texttt{Totale waarde:} \\
\texttt{Jij: 2 bloemen = 8} \\
\texttt{Ik: 2 tangen + 2 stokken = 6 + 2 = 8} \\
\texttt{Totaal: 16} \\
\\ 
\texttt{Als we de bloemen aan jou geven en de tangen aan mij, en de stokken willekeurig verdelen, is de totale waarde maximaal.} \\
\\ 
\texttt{Zou jij akkoord gaan met de volgende verdeling:} \\
\texttt{{-} Jij: 2 bloemen} \\
\texttt{{-} Ik: 2 tangen, 2 stokken} \\
\\ 
\texttt{Dit geeft jou 8 punten en mij 8 punten, totaal 16.} \\
\\ 
\texttt{Laat me weten of dit voorstel voor jou acceptabel is, dan kunnen we het geheim voorstel indienen.} \\
            }
        }
    }
    & & \\ \\

    \theutterance \stepcounter{utterance}  
    & & \multicolumn{4}{p{0.6\linewidth}}{
        \cellcolor[rgb]{0.9,0.9,0.9}{
            \makecell[{{p{\linewidth}}}]{
                \texttt{\tiny{[P1$\langle$GM]}}
                \texttt{Ik waardeer de items als volgt:} \\
\texttt{{-} Bloem: 1} \\
\texttt{{-} Tang: 3} \\
\texttt{{-} Stok: 1} \\
\\ 
\texttt{Op basis van jouw waarden (bloem: 4, tang: 1, stok: 0) en de beschikbare items (2 bloemen, 2 tangen, 2 stokken), lijkt de optimale verdeling voor de totale som als volgt:} \\
\\ 
\texttt{Jij krijgt de 2 bloemen (waarde voor jou: 4 + 4 = 8).} \\
\texttt{Ik krijg de 2 tangen (waarde voor mij: 3 + 3 = 6).} \\
\texttt{De 2 stokken kunnen we verdeeld houden of aan elkaar geven, maar aangezien jij ze op 0 waardeert en ik op 1, is het voor de totale som neutraal of zelfs beter als ik ze krijg. Laten we zeggen dat jij de bloemen krijgt en ik de tangen en stokken.} \\
\\ 
\texttt{Totale waarde:} \\
\texttt{Jij: 2 bloemen = 8} \\
\texttt{Ik: 2 tangen + 2 stokken = 6 + 2 = 8} \\
\texttt{Totaal: 16} \\
\\ 
\texttt{Als we de bloemen aan jou geven en de tangen aan mij, en de stokken willekeurig verdelen, is de totale waarde maximaal.} \\
\\ 
\texttt{Zou jij akkoord gaan met de volgende verdeling:} \\
\texttt{{-} Jij: 2 bloemen} \\
\texttt{{-} Ik: 2 tangen, 2 stokken} \\
\\ 
\texttt{Dit geeft jou 8 punten en mij 8 punten, totaal 16.} \\
\\ 
\texttt{Laat me weten of dit voorstel voor jou acceptabel is, dan kunnen we het geheim voorstel indienen.} \\
            }
        }
    }
    & \\ \\

    \theutterance \stepcounter{utterance}  
    & \multicolumn{4}{p{0.6\linewidth}}{
        \cellcolor[rgb]{0.9,0.9,0.9}{
            \makecell[{{p{\linewidth}}}]{
                \texttt{\tiny{[P1$\rangle$GM]}}
                \texttt{Dat klinkt als een uitstekend plan! Door de bloemen aan mij te geven en de tangen (en stokken) aan jou, maximaliseren we inderdaad de totale waarde.} \\
\\ 
\texttt{Ik ga akkoord met deze verdeling. Laten we het geheim voorstel indienen.} \\
\\ 
\texttt{{[}Voorstel: 2 bloemen{]}} \\
            }
        }
    }
    & & \\ \\

    \theutterance \stepcounter{utterance}  
    & & & \multicolumn{2}{p{0.3\linewidth}}{
        \cellcolor[rgb]{0.9,0.9,0.9}{
            \makecell[{{p{\linewidth}}}]{
                \texttt{\tiny{[GM$|$GM]}}
                \texttt{proposal} \\
            }
        }
    }
    & & \\ \\

    \theutterance \stepcounter{utterance}  
    & & & \multicolumn{2}{p{0.3\linewidth}}{
        \cellcolor[rgb]{0.9,0.9,0.9}{
            \makecell[{{p{\linewidth}}}]{
                \texttt{\tiny{[GM$|$GM]}}
                [2, 0, 0]
            }
        }
    }
    & & \\ \\

    \theutterance \stepcounter{utterance}  
    & & \multicolumn{4}{p{0.6\linewidth}}{
        \cellcolor[rgb]{0.9,0.9,0.9}{
            \makecell[{{p{\linewidth}}}]{
                \texttt{\tiny{[GM$\rangle$P2]}}
                \texttt{De andere speler heeft een geheim voorstel ingediend. Het is nu jouw beurt om een eigen voorstel in te dienen.} \\
            }
        }
    }
    & \\ \\

    \theutterance \stepcounter{utterance}  
    & & & \multicolumn{4}{p{0.6\linewidth}}{
        \cellcolor[rgb]{0.9,0.9,0.9}{
            \makecell[{{p{\linewidth}}}]{
                \texttt{\tiny{[GM$\langle$P2]}}
                \texttt{{[}Forstel: 2 bloemen, 2 stokken{]}} \\
            }
        }
    }
     \\ \\

    \theutterance \stepcounter{utterance}  
    & & & \multicolumn{2}{p{0.3\linewidth}}{
        \cellcolor[rgb]{0.9,0.9,0.9}{
            \makecell[{{p{\linewidth}}}]{
                \texttt{\tiny{[GM$|$GM]}}
                \texttt{proposal does not start with "voorstel:"} \\
            }
        }
    }
    & & \\ \\

\end{supertabular}
}

\end{document}
